% !PRINTED 585
\weberpage{585}

It is then, by (2),
\weq{8}{\begin{aligned}
\mu &= y + bx + 2a\omega x, \\
\mu\omega &= -2cx + (y-bx)\omega, \\
m &= (y+bx)(y-bx) + 4acx^2,
\end{aligned}}
from which it follows that
\[
y + bx \text{ and } y - bx
\]
are odd.

From this there results, by (4):
\weq{9}{\begin{aligned}
\boldsymbol{s}[\mu(u+1)] &= -\boldsymbol{s}(\mu u), \\
\boldsymbol{s}[\mu(u+\omega)] &= \boldsymbol{s}(\mu u).
\end{aligned}}
The function $\boldsymbol{s}(\mu u)$ therefore has the same periods as $\boldsymbol{s}(u)$,
and since $\boldsymbol{s}(\mu u)$ is an odd function of $u$, it can be expressed rationally by $\boldsymbol{s}(u)$. We denote by $A(\boldsymbol{s})$, $D(\boldsymbol{s})$ integral rational functions of $\boldsymbol{s} = \boldsymbol{s}(u)$, which moreover contain only the even powers of $\boldsymbol{s}$, and set:
\weq{10}{\boldsymbol{s}(\mu u) = \mu\,\boldsymbol{s}(u)\,\dfrac{A(\boldsymbol{s})}{D(\boldsymbol{s})}.}
The values of $\boldsymbol{s}(u)$ for which $\boldsymbol{s}(\mu u)$ becomes zero or infinite are, by 3., algebraic numbers, and since the coefficients of $A$ and $D$ can be composed out of these, these too are algebraic numbers. In order to determine their nature more closely, we think of the fraction (10) as first extended in such a way that $D$ receives rational coefficients. This is brought about by multiplying numerator and denominator by the product of all the factors conjugate to the denominator. We can therefore put $A(\boldsymbol{s})$ and $D(\boldsymbol{s})$ into the form:
\weq{11}{\begin{aligned}
A(\boldsymbol{s}) &= A_1 + A_3\boldsymbol{s}^2 + A_5\boldsymbol{s}^4 + \cdots, \\
D(\boldsymbol{s}) &= D_1 + D_3\boldsymbol{s}^2 + D_5\boldsymbol{s}^4 + \cdots,
\end{aligned}}
in which the $D_1, D_3, \dots$, are by our assumption rational numbers, while $A_1, A_3, \dots$, are to be determined. If we set, as in \S\,155, (4), (6):
\[
w = \sqrt{\varkappa}\,v = 2\sqrt{\varkappa}\,Ku,
\]
then $\boldsymbol{s}(u)$ goes over into $\sqrt{\varkappa}\ \mathrm{sn}\left(\dfrac{w}{\sqrt{\varkappa}}\right)$, and equation (10) gives:
\weq{12}{\dfrac{D\left[\sqrt{\varkappa}\ \mathrm{sn}\left(\dfrac{\omega}{\sqrt{\varkappa}}\right)\right]\ \sqrt{\varkappa}\ \mathrm{sn}\left(\dfrac{\mu w}{\sqrt{\varkappa}}\right)}{\mu\ \mathrm{sn}\left(\dfrac{w}{\sqrt{\varkappa}}\right)} = A_1 + A_3\boldsymbol{s}^2 + A_5\boldsymbol{s}^4 + \cdots}
% [source sic: D's argument prints \omega while the other arguments use w, and the denominator lacks the otherwise matching \sqrt{\varkappa}]
